% ================================================================
% PROCESSI ORGANIZZATIVI
% ================================================================

\section{Processi Organizzativi}
\label{sec:processi-organizzativi}

In conformità allo standard ISO/IEC~12207:1995, i processi
organizzativi comprendono le attività trasversali necessarie a governare il gruppo,
mantenere l'infrastruttura di lavoro, favorire il miglioramento continuo e garantire
l'acquisizione delle competenze necessarie.

Nel contesto del progetto \textit{Second Brain}, e in accordo con il tailoring
descritto nella Sezione~\ref{sec:tailoring}, il gruppo \textit{Synergo Unit} adotta
i seguenti processi organizzativi:

\begin{itemize}
    \item gestione;
    \item gestione dei rischi;
    \item comunicazione;
    \item infrastruttura;
    \item miglioramento;
    \item formazione.
\end{itemize}

% ================================================================
% GESTIONE
% ================================================================

\subsection{Processo di Gestione}
\label{sec:gestione}

\subsubsection{Descrizione}

Il processo di gestione disciplina l'organizzazione interna del gruppo, la
pianificazione delle attività, l'assegnazione e la rotazione dei ruoli, la
conduzione degli sprint, il tracciamento dell'avanzamento e il
monitoraggio di tempi e costi.

Esso fornisce la struttura organizzativa entro cui i processi primari e di supporto
vengono attuati.

% ------------------------------------------------

\subsubsection{Composizione del Gruppo}

Il gruppo \textit{Synergo Unit} è composto da sei membri.

Ogni membro contribuisce allo svolgimento del progetto assumendo, nel tempo, i
ruoli previsti dal Regolamento del Progetto Didattico.

La composizione del gruppo, la rotazione dei ruoli e la distribuzione delle ore
sono riportate nel Piano di Progetto.

% ------------------------------------------------

\subsubsection{Ruoli del Gruppo}
\label{sec:ruoli}

Il gruppo adotta i seguenti ruoli:

\begin{itemize}
    \item Responsabile;
    \item Amministratore;
    \item Analista;
    \item Progettista;
    \item Programmatore\textsubscript{G};
    \item Verificatore.
\end{itemize}

Ogni membro può ricoprire un solo ruolo all'interno dello stesso sprint. Tuttavia,
uno stesso ruolo può essere assegnato a più membri nello stesso sprint, qualora il
carico di lavoro lo renda necessario.

In fase RTB, il ruolo di Programmatore è stato attivato solo per attività di codifica
limitate alla realizzazione del Proof of Concept. La piena attivazione del ruolo è avvenuta
con l'avvio dello sviluppo del Minimum Viable Product nella fase PB.

La rotazione dei ruoli è definita sprint per sprint durante la riunione di
pianificazione. Essa tiene conto di:

\begin{itemize}
    \item disponibilità temporale dei membri;
    \item carico di lavoro previsto;
    \item competenze richieste dalle attività;
    \item necessità di distribuire l'esperienza sui diversi ruoli;
    \item obiettivo di mantenere tutti i membri aggiornati sull'avanzamento del progetto.
\end{itemize}

Le assegnazioni sono decise collegialmente dal gruppo durante il planning e
registrate nel verbale interno.

La rotazione dei ruoli e l'assegnazione delle attività devono favorire la
distribuzione delle competenze all'interno del gruppo.

Il gruppo non adotta una logica di assegnazione basata esclusivamente sulla
competenza pregressa del singolo membro: la scelta dell'assegnatario non deve
ricadere automaticamente su chi possiede già maggiore familiarità con una
specifica attività o tecnologia.

Quando una competenza è concentrata in uno o pochi membri, il gruppo deve
favorire affiancamento, condivisione delle conoscenze e progressiva autonomia
degli altri componenti, compatibilmente con scadenze, carico di lavoro e criticità
dello sprint.

% ------------------------------------------------

\subsubsection{Costi Orari dei Ruoli}

I costi orari adottati sono quelli previsti dal Regolamento del Progetto Didattico.

\begin{xltabular}{\textwidth}{
    >{\raggedright\arraybackslash}p{3.5cm}
    >{\raggedright\arraybackslash}p{3cm}
    >{\raggedright\arraybackslash}X
}
    \toprule
    \textbf{Ruolo} & \textbf{Costo Orario} & \textbf{Responsabilità Principale} \\
    \midrule
    \endfirsthead

    \toprule
    \textbf{Ruolo} & \textbf{Costo Orario} & \textbf{Responsabilità Principale} \\
    \midrule
    \endhead

    \midrule
    \multicolumn{3}{r}{\textit{Continua nella pagina successiva}} \\
    \endfoot

    \bottomrule
    \endlastfoot

    Responsabile
    & 30\,\euro{}/h
    & Coordina il gruppo, rappresenta il fornitore verso committente e proponente,
    approva il rilascio degli artefatti e supervisiona l'avanzamento complessivo. \\

    Amministratore
    & 20\,\euro{}/h
    & Assicura l'efficienza degli strumenti e delle procedure a supporto del
    way of working, mantiene le Norme di Progetto e gestisce
    l'infrastruttura operativa. \\

    Analista
    & 25\,\euro{}/h
    & Svolge attività di analisi dei requisiti, studio del dominio e confronto tra
    esigenze della proponente e funzionalità attese. \\

    Progettista
    & 25\,\euro{}/h
    & Svolge attività di progettazione, selezione tecnologica e definizione della
    soluzione tecnica. \\

    Programmatore
    & 15\,\euro{}/h
    & Svolge attività di codifica. In fase PB implementa le funzionalità del Minimum 
    Viable Product, scrive i test automatici e assicura che il codice rispetti le 
    metriche di analisi statica. \\

    Verificatore
    & 15\,\euro{}/h
    & Svolge attività di verifica sugli artefatti prodotti dagli altri membri,
    garantendo indipendenza rispetto al Redattore. \\

\end{xltabular}

% ------------------------------------------------

\subsubsection{Responsabilità Operative Specifiche}

Oltre alle responsabilità previste dal ruolo, il gruppo adotta le seguenti
responsabilità operative:

\begin{itemize}
    \item il \textbf{Responsabile} coordina le riunioni, approva i documenti
    verificati, supervisiona l'avanzamento complessivo e cura il coordinamento
    verso committente e proponente;

    \item l'\textbf{Amministratore} mantiene l'infrastruttura, gestisce Google
    Forms, Google Sheets, GitHub Pages,
    email istituzionale, verbali esterni e aggiornamenti delle Norme di Progetto;

    \item l'\textbf{Analista} cura l'Analisi dei Requisiti, le fonti, i casi d'uso,
    il tracciamento e il confronto con la proponente sugli aspetti funzionali;

    \item il \textbf{Progettista} cura la valutazione tecnica delle soluzioni,
    lo stack tecnologico e il perimetro tecnico del Proof of Concept;

    \item il \textbf{Programmatore} realizza codice solo quando necessario per il
    Proof of Concept in RTB e in modo più esteso nella fase PB;

    \item il \textbf{Verificatore} controlla gli artefatti prodotti da altri membri,
    senza modificare direttamente il lavoro verificato.
\end{itemize}

% ------------------------------------------------

\subsubsection{Pianificazione degli Sprint}

Gli sprint hanno durata settimanale, con cadenza da martedì a martedì.

Le attività vengono pianificate durante la riunione del martedì, dopo la
retrospettiva e la review dello sprint precedente.

Durante lo sprint planning il gruppo:

\begin{itemize}
    \item individua le attività da svolgere;
    \item valuta priorità e scadenze;
    \item stima i tempi previsti;
    \item assegna i ruoli dello sprint;
    \item assegna i task ai membri;
    \item crea le issue tramite Google Form;
    \item aggiorna il backlog\textsubscript{G} dello sprint.
\end{itemize}

I task vengono pianificati ordinariamente durante il planning del martedì.
L'aggiunta di nuovi task durante lo sprint è ammessa solo in caso di necessità
motivata e deve essere tracciata.

% ------------------------------------------------

\subsubsection{Gestione dei Task Non Completati}

I task non completati entro la fine dello sprint non vengono chiusi e ricreati.

Essi devono essere:

\begin{itemize}
    \item discussi durante la sprint review;
    \item motivati nel verbale interno;
    \item mantenuti aperti;
    \item riportati nella lista delle issue dello sprint in cui vengono posticipati.
\end{itemize}

Per i task non associati a modifiche della repository, la chiusura della issue
spetta all'assegnatario, dopo il completamento dell'attività e la verifica in
review.

% ------------------------------------------------

\subsubsection{Tracciamento Ore e Costi}
\label{sec:calcolo-ore}

Il tracciamento delle ore avviene in minuti.

Il tempo previsto viene inserito nel Google Form al momento della creazione del
task. Il tempo effettivo viene inserito manualmente nel foglio Google condiviso:

\begin{itemize}
    \item dopo il merge e la chiusura della issue, per attività con modifica della
    repository;
    \item al termine dell'attività, per task non associati a modifiche della repository.
\end{itemize}

Per ogni task devono essere registrati:

\begin{itemize}
    \item tempo previsto del ruolo assegnatario;
    \item tempo effettivo del ruolo assegnatario;
    \item tempo previsto ed effettivo del Verificatore;
    \item tempo previsto ed effettivo del Responsabile, se coinvolto;
    \item tempo previsto ed effettivo dell'Amministratore, se coinvolto.
\end{itemize}

L'Amministratore controlla che i tempi effettivi siano compilati correttamente e
li utilizza per:

\begin{itemize}
    \item consuntivo\textsubscript{G} di sprint;
    \item aggiornamento del Piano di Progetto;
    \item monitoraggio dei costi;
    \item analisi delle risorse residue\textsubscript{G}.
\end{itemize}

Il monitoraggio dei costi deve garantire il rispetto del costo massimo accettato
in fase di aggiudicazione. Qualora i costi sostenuti si avvicinino al limite
previsto, il gruppo deve valutare azioni correttive e, se necessario, rinegoziare
il perimetro dei requisiti con la proponente.

Il foglio di calcolo\textsubscript{G} condiviso produce automaticamente:

\begin{itemize}
    \item preventivo di sprint;
    \item consuntivo di sprint;
    \item delta ore;
    \item delta costi;
    \item risorse residue per ruolo;
    \item risorse residue per membro.
\end{itemize}

Le formule di conversione applicate sono:

\begin{align*}
  \text{ore} &= \frac{\text{minuti}}{60} \\
  \text{costo} &= \frac{\text{minuti}}{60} \times \text{costo orario del ruolo}
\end{align*}

\subsubsection{Gestione dei Rischi}

La gestione dei rischi è responsabilità del Responsabile e viene formalizzata nel
Piano di Progetto.

Durante ogni sprint retrospective, il gruppo valuta:
\begin{itemize}
    \item rischi occorsi;
    \item variazioni di probabilità o impatto dei rischi già identificati;
    \item nuovi rischi emersi;
    \item efficacia delle strategie di mitigazione\textsubscript{G} adottate.
\end{itemize}

Gli aggiornamenti rilevanti vengono riportati nel Piano di Progetto e, se
necessario, trasformati in azioni correttive o task operativi.

% ------------------------------------------------

\subsubsection{Riunione del Martedì}

La riunione ordinaria del martedì si tiene alle 14:30 su Discord\textsubscript{G}.

Qualora emergano impedimenti, l'orario può essere anticipato o posticipato previo
accordo su WhatsApp\textsubscript{G}. L'orario effettivo di inizio e fine viene
riportato nel verbale interno.

La riunione ha durata variabile in base
alla complessità delle decisioni da assumere.

La struttura della riunione è:

\begin{enumerate}
    \item \textbf{Sprint Review}: il gruppo verifica lo stato dei task, analizza
    eventuali scostamenti tra preventivo e consuntivo e valuta i task non completati;

    \item \textbf{Sprint Retrospective}: il gruppo analizza il processo seguito nello
    sprint precedente e individua almeno un'azione correttiva;

    \item \textbf{Sprint Planning}: il gruppo pianifica lo sprint corrente, assegna
    ruoli e task e crea le relative issue.
\end{enumerate}

Le decisioni, le azioni e le assegnazioni emerse vengono registrate nel documento
di riunione condiviso e successivamente formalizzate nel verbale interno.

% ------------------------------------------------

\subsubsection{Riunione del Venerdì}

La riunione del venerdì è fissata alle 14:30 su Discord e ha durata indicativa di
30 minuti.

La riunione può essere anticipata o posticipata, di comune accordo, al termine del
Diario di Bordo, in modo da recepire tempestivamente eventuali aggiornamenti o
osservazioni.

Essa ha lo scopo di:

\begin{itemize}
    \item risolvere dubbi emersi durante lo sprint;
    \item affrontare criticità operative;
    \item aggiornare il gruppo sullo stato delle attività;
    \item recepire eventuali osservazioni emerse dal Diario di Bordo.
\end{itemize}

% ------------------------------------------------

\subsubsection{Daily Stand-Up Asincrono}

Il gruppo utilizza WhatsApp per comunicazioni operative e coordinamento durante
lo sprint.

Il Daily Stand-Up asincrono viene utilizzato quando utile per segnalare
impedimenti, richieste di supporto, dubbi o aggiornamenti rilevanti
rispetto alle attività in corso.

Lo stato di avanzamento delle attività non viene riportato sistematicamente
tramite il Daily Stand-Up, poiché è tracciato in modo continuativo mediante
GitHub Projects\textsubscript{G}, issue\textsubscript{G}, branch\textsubscript{G}
e pull request\textsubscript{G}.

Quando necessario, i membri del gruppo possono comunque comunicare:

\begin{itemize}
    \item avanzamenti significativi;
    \item attività previste nel breve termine;
    \item impedimenti;
    \item richieste di supporto.
\end{itemize}

% ------------------------------------------------

\subsubsection{Verbali e Tracciamento Decisioni}

Ogni riunione interna produce un verbale interno.

Il verbale interno deve riportare:

\begin{itemize}
    \item data, orario e piattaforma;
    \item partecipanti presenti e assenti;
    \item ordine del giorno;
    \item sintesi della discussione;
    \item decisioni prese;
    \item azioni da svolgere;
    \item riferimenti alle issue create;
    \item data della riunione successiva.
\end{itemize}

Le decisioni interne sono identificate tramite codice:

\begin{center}
    \texttt{VI-y.x}
\end{center}

dove \texttt{y} identifica il verbale e \texttt{x} la decisione progressiva.

Ogni decisione rilevante deve essere collegata, quando possibile, a una o più
issue operative.

% ================================================================
% COMUNICAZIONE
% ================================================================

\subsection{Processo di Comunicazione}
\label{sec:comunicazione}

\textit{Nota di tailoring}: lo standard ISO/IEC~12207:1995 definisce quattro
processi organizzativi (Gestione, Infrastruttura, Miglioramento, Formazione).
Il presente processo non corrisponde a nessuno dei quattro previsti dallo
standard: è un'estensione deliberata, introdotta dal gruppo per normare
esplicitamente canali e responsabilità di comunicazione che lo standard non
tratta come processo a sé, ritenuta necessaria per la natura distribuita e
la dimensione del gruppo (Sezione~\ref{sec:tailoring}).

\subsubsection{Descrizione}

Il processo di comunicazione disciplina i canali, le modalità e le responsabilità
relative alle comunicazioni interne ed esterne del gruppo.

Una comunicazione corretta deve essere:

\begin{itemize}
    \item tracciabile;
    \item tempestiva;
    \item pertinente;
    \item accessibile ai membri coinvolti;
    \item formalizzata quando produce decisioni o impegni.
\end{itemize}

% ------------------------------------------------

\subsubsection{Comunicazione Interna}

I canali interni adottati sono:

\begin{itemize}
    \item \textbf{Discord}: piattaforma principale per riunioni, discussioni operative
    e allineamenti sincroni;

    \item \textbf{WhatsApp}: canale per comunicazioni rapide, aggiornamenti asincroni,
    impedimenti e convocazioni straordinarie;

    \item \textbf{GitHub}: canale operativo per issue, Pull Request, review e
    tracciamento delle attività;

    \item \textbf{Google Calendar\textsubscript{G}}: strumento per disponibilità, eventi ricorrenti e
    pianificazione delle riunioni;

    \item \textbf{Google Drive/Docs}: spazio di supporto per documenti di riunione
    condivisi e materiali temporanei.
\end{itemize}

Le decisioni prese su canali informali devono essere riportate nel primo verbale
interno utile, qualora producano effetti su documenti, task, requisiti, strumenti
o processi.

% ------------------------------------------------

\subsubsection{Comunicazione Esterna}

Le comunicazioni esterne avvengono tramite l'email istituzionale del gruppo:

\begin{center}
    \texttt{synergounit.20@gmail.com}
\end{center}

L'indirizzo è gestito dall'Amministratore.

L'Amministratore cura la gestione operativa della casella email; il Responsabile
mantiene la responsabilità del contenuto delle comunicazioni formali quando esse
comportano impegni verso committente o proponente.

Le comunicazioni esterne devono passare attraverso l'email istituzionale, salvo
diversa indicazione esplicita della proponente o del committente.

% ------------------------------------------------

\subsubsection{Incontri con la Proponente}

Gli incontri con la proponente avvengono tramite Google Meet.

Il processo prevede:

\begin{enumerate}
    \item durante una riunione interna il gruppo individua i temi da discutere;
    \item il gruppo definisce l'ordine del giorno;
    \item l'Amministratore contatta la proponente tramite email ufficiale;
    \item dopo conferma della disponibilità, il gruppo invia l'evento Google Meet;
    \item prima dell'incontro vengono definiti i portavoce per i diversi argomenti;
    \item al termine dell'incontro l'Amministratore redige il verbale esterno;
    \item il verbale viene inviato alla proponente per firma;
    \item dopo conferma, il verbale viene caricato, verificato, approvato e pubblicato.
\end{enumerate}

Il portavoce dell'incontro viene scelto in base all'argomento trattato. In caso di
più tematiche, il gruppo stabilisce preventivamente chi interviene su ciascun
punto, mantenendo comunque la possibilità per ogni membro di porre domande.

L'obiettivo è ottimizzare il tempo dell'incontro, evitare tempi morti e garantire
che tutti i dubbi preparati vengano affrontati.

% ------------------------------------------------

\subsubsection{Comunicazione con il Committente}

La comunicazione con il committente avviene tramite:

\begin{itemize}
    \item email istituzionale del gruppo;
    \item Diari di Bordo;
    \item revisioni di avanzamento;
    \item incontri o ricevimenti richiesti.
\end{itemize}

Gli esiti rilevanti delle comunicazioni con il committente devono essere analizzati
dal gruppo e, se necessario, tradotti in azioni operative.

% ================================================================
% INFRASTRUTTURA
% ================================================================

\subsection{Processo di Infrastruttura}
\label{sec:infrastruttura}

\subsubsection{Descrizione}

Il processo di infrastruttura definisce e mantiene l'insieme degli strumenti
utilizzati dal gruppo per supportare pianificazione, comunicazione, documentazione,
sviluppo, pubblicazione e tracciamento.

L'infrastruttura deve essere mantenuta coerente con i processi definiti nelle
presenti Norme di Progetto.

% ------------------------------------------------

\subsubsection{Responsabilità}

L'Amministratore è responsabile della manutenzione dell'infrastruttura operativa.

In particolare, l'Amministratore:

\begin{itemize}
    \item mantiene configurati gli strumenti di gestione;
    \item mantiene Google Forms e Google Sheets;
    \item gestisce l'email istituzionale;
    \item aggiorna il sito GitHub Pages;
    \item supporta i membri nell'utilizzo degli strumenti;
    \item propone aggiornamenti alle Norme di Progetto quando cambia l'infrastruttura.
\end{itemize}

% ------------------------------------------------

\subsubsection{Strumenti di Gestione}

Gli strumenti di gestione adottati sono:

\begin{itemize}
    \item \textbf{GitHub Projects}: pianificazione e tracciamento dello stato delle
    issue;

    \item \textbf{Google Forms}: creazione standardizzata dei task e generazione
    automatica delle issue;

    \item \textbf{Google Sheets}: raccolta di stime, consuntivi e calcolo automatico
    di ore, costi e risorse residue;

    \item \textbf{Google Calendar}: gestione di riunioni, disponibilità ed eventi.
\end{itemize}

% ------------------------------------------------

\subsubsection{Strumenti di Documentazione}

Gli strumenti di documentazione adottati sono:

\begin{itemize}
    \item \textbf{\LaTeX{}};
    \item \textbf{Visual Studio Code};
    \item \textbf{LaTeX Workshop};
    \item \textbf{Script Python per Glossario};
    \item \textbf{GitHub Pages};
    \item \textbf{Canva}.
\end{itemize}

Le modalità operative di utilizzo di tali strumenti sono descritte nel Processo
di Documentazione e nel Processo di Gestione della Configurazione.

% ------------------------------------------------

\subsubsection{Strumenti di Comunicazione}

Gli strumenti di comunicazione adottati sono:

\begin{itemize}
    \item \textbf{Discord};
    \item \textbf{WhatsApp};
    \item \textbf{Email istituzionale};
    \item \textbf{Google Meet};
    \item \textbf{Google Drive/Docs}.
\end{itemize}

% ------------------------------------------------

\subsubsection{Strumenti Tecnici}

Gli strumenti e le tecnologie tecniche adottati per il Proof of Concept includono:

\begin{itemize}
    \item Vue.js;
    \item CodeMirror;
    \item Python;
    \item FastAPI;
    \item Docker;
    \item Docker Compose;
    \item GitHub;
    \item LucidSpark con elementi UML di LucidChart.
    \item StarUML
\end{itemize}

In fase RTB il Proof of Concept non prevede persistenza su database. L'eventuale
adozione di PostgreSQL è stata rivalutata nella fase PB e, a valle del
confronto con la proponente, il gruppo ha deciso di non implementare la
persistenza server-side in questa release (R4-V-Op, non implementato):
la gestione dei dati si basa esclusivamente sul file system locale del
browser, per concentrare le risorse disponibili sul consolidamento dei
requisiti obbligatori (dettaglio in Specifica Tecnica, Sezione ``Persistenza
dei dati'').

Le decisioni definitive relative allo stack tecnologico sono riportate
nell'Analisi dello Stato dell'Arte e potranno essere aggiornate in seguito al
confronto con la proponente e agli esiti del Proof of Concept.

% ------------------------------------------------

\subsubsection{Infrastruttura del Proof of Concept}

L'infrastruttura del Proof of Concept è contenuta nella cartella \texttt{poc/}
della repository \texttt{Second-Brain}.

Il PoC viene eseguito tramite Docker Compose e comprende due servizi:

\begin{itemize}
    \item \texttt{frontend}: servizio dedicato all'interfaccia Vue.js;
    \item \texttt{backend}: servizio dedicato alle API FastAPI.
\end{itemize}

La procedura di avvio del PoC deve essere documentata nel README della cartella
\texttt{poc/}.

Il README deve contenere almeno:

\begin{itemize}
    \item prerequisiti;
    \item procedura di avvio tramite Docker Compose;
    \item porte utilizzate dai servizi;
    \item modalità di verifica dell'avvio corretto;
    \item eventuali problemi noti.
\end{itemize}

Il README tecnico del PoC è mantenuto dal Progettista quando riguarda scelte
tecniche o architetturali, e dal Programmatore quando riguarda comandi operativi,
setup ed esecuzione locale.

% ------------------------------------------------

\subsubsection{Struttura del Repository di Sviluppo}

In fase PB, il codice sorgente è organizzato all'interno del repository in due macro-cartelle principali, al fine di garantire una netta separazione delle responsabilità tra i due domini.

\paragraph{Frontend (Vue.js / TypeScript)}
L'applicazione lato client è organizzata secondo il pattern MVC. La struttura interna dei package, situata all'interno della cartella \texttt{src/}, è la seguente:
\begin{itemize}
    \item \textbf{\texttt{model/}}: contiene lo stato applicativo e la logica dei dati;
    \item \textbf{\texttt{view/}}: contiene i componenti grafici e di presentazione Vue;
    \item \textbf{\texttt{controller/}}: gestisce e smista gli eventi generati dall'utente;
    \item \textbf{\texttt{proxy/}}: implementa le interfacce di servizio traducendo le chiamate verso il backend;
    \item \textbf{\texttt{bootstrap/}}: gestisce la composizione manuale e l'iniezione delle dipendenze.
\end{itemize}

\paragraph{Backend (FastAPI / Python)}
L'applicazione lato server adotta una \textit{Layered Architecture}. La struttura interna dei package, situata all'interno della cartella \texttt{app/}, è la seguente:
\begin{itemize}
    \item \textbf{\texttt{api/}}: \textit{Presentation Layer} per l'esposizione degli endpoint REST e la validazione;
    \item \textbf{\texttt{service/}}: \textit{Business Layer} per l'orchestrazione dei casi d'uso;
    \item \textbf{\texttt{domain/}}: contiene le entità di dominio, i modelli e le operazioni AI;
    \item \textbf{\texttt{llm/}}: contiene le interfacce e gli adapter per comunicare con i provider LLM;
    \item \textbf{\texttt{export/}}: gestisce le implementazioni per i vari formati di esportazione;
    \item \textbf{\texttt{di/}}: contiene i provider per la Dependency Injection.
\end{itemize}

% ------------------------------------------------

\subsubsection{Adozione di Nuovi Strumenti}

L'adozione di un nuovo strumento deve seguire la seguente procedura:

\begin{enumerate}
    \item proposta durante una riunione interna;
    \item valutazione collegiale di utilità, costi, benefici e impatto sul workflow;
    \item approvazione del gruppo;
    \item registrazione della decisione nel verbale interno;
    \item aggiornamento delle Norme di Progetto prima dell'utilizzo operativo.
\end{enumerate}

% ================================================================
% MIGLIORAMENTO
% ================================================================

\subsection{Processo di Miglioramento}
\label{sec:miglioramento}

\subsubsection{Descrizione}

Il processo di miglioramento ha lo scopo di individuare inefficienze nei processi
attivi e introdurre azioni correttive finalizzate al miglioramento continuo del
way of working.

Il processo è collegato al ciclo PDCA descritto nella Sezione~\ref{sec:qualita}
e viene attuato principalmente tramite retrospective, consuntivi di sprint e
aggiornamento delle presenti Norme di Progetto.

% ------------------------------------------------

\subsubsection{Identificazione delle Criticità}

Le criticità possono emergere da:

\begin{itemize}
    \item sprint review;
    \item sprint retrospective;
    \item scostamenti tra preventivo e consuntivo;
    \item difficoltà operative riportate dai membri;
    \item non conformità;
    \item feedback di committente o proponente;
    \item problemi nell'utilizzo degli strumenti.
\end{itemize}

% ------------------------------------------------

\subsubsection{Azioni Correttive}

Ogni sprint retrospective deve produrre almeno una azione correttiva, anche minima.

Le azioni correttive devono essere:

\begin{itemize}
    \item verbalizzate nel verbale interno;
    \item riportate nel consuntivo di sprint;
    \item monitorate nello sprint successivo;
    \item trasformate in aggiornamenti delle Norme di Progetto, se modificano il
    way of working.
\end{itemize}

Le azioni correttive possono riguardare:

\begin{itemize}
    \item processo;
    \item comunicazione;
    \item strumenti;
    \item pianificazione;
    \item qualità documentale;
    \item gestione dei task;
    \item stima e consuntivazione.
\end{itemize}

% ------------------------------------------------

\subsubsection{Aggiornamento del Way of Working}

Nessun processo può essere modificato unilateralmente.

Ogni modifica al way of working deve:

\begin{itemize}
    \item essere discussa in riunione;
    \item essere approvata dal gruppo;
    \item essere verbalizzata;
    \item essere recepita nelle Norme di Progetto dall'Amministratore.
\end{itemize}

% ================================================================
% FORMAZIONE
% ================================================================

\subsection{Processo di Formazione}
\label{sec:formazione}

\subsubsection{Descrizione}

Il processo di formazione ha lo scopo di garantire che i membri del gruppo
acquisiscano le competenze necessarie per utilizzare strumenti, tecnologie e
processi adottati.

La formazione riduce i rischi derivanti da inesperienza tecnica, rotazione dei
ruoli e introduzione di nuovi strumenti.

% ------------------------------------------------

\subsubsection{Modalità di Formazione}

La formazione viene pianificata tramite task specifici quando una nuova tecnologia,
uno strumento o una procedura richiede studio preliminare.

Le attività di formazione possono avvenire tramite:

\begin{itemize}
    \item studio individuale della documentazione ufficiale;
    \item prove pratiche su repository o ambienti locali;
    \item confronto tra membri;
    \item sessioni di allineamento su Discord;
    \item condivisione di indicazioni operative da parte di membri più esperti.
\end{itemize}

Le ore dedicate alla formazione sono considerate ore di palestra e non vengono
consuntivate come ore produttive di progetto.

% ------------------------------------------------

\subsubsection{Formazione su Strumenti di Processo}

La formazione sugli strumenti di processo è coordinata dall'Amministratore.

Essa riguarda principalmente:

\begin{itemize}
    \item GitHub Projects;
    \item Google Forms;
    \item Google Sheets;
    \item GitHub Pages;
    \item template documentali;
    \item script di automazione del Glossario.
\end{itemize}

% ------------------------------------------------

\subsubsection{Formazione Tecnica}

La formazione tecnica è coordinata dai membri che ricoprono i ruoli di Analista e
Progettista, in base all'ambito trattato.

Essa riguarda principalmente:

\begin{itemize}
    \item Vue.js;
    \item CodeMirror;
    \item Python;
    \item FastAPI;
    \item PostgreSQL, qualora venga confermato nelle successive decisioni progettuali;
    \item Docker;
    \item strumenti UML;
    \item strumenti di sviluppo collegati al Proof of Concept.
\end{itemize}

Le risorse di riferimento principali sono le documentazioni ufficiali delle
tecnologie adottate.

% ------------------------------------------------

\subsubsection{Applicazione della Formazione}

La formazione deve essere avviata prima dell'utilizzo operativo dello strumento o
della tecnologia.

Qualora l'apprendimento possa avvenire in modo incrementale, il membro può
proseguire la formazione durante l'attività, purché siano garantiti:

\begin{itemize}
    \item comprensione minima dello strumento;
    \item supporto da parte del gruppo;
    \item verifica dell'artefatto prodotto;
    \item tracciabilità dell'attività svolta.
\end{itemize}
